\section{Introduction}\label{sec:intro}

Smishing, or phishing by short text message, imitates a bank, a delivery service or a government agency and leads the recipient to a link, a phone number or an application in order to obtain credentials or a payment. Unlike advertising spam, a missed smishing message can cause direct financial loss, so the two kinds of error have unequal costs. A detector for this task works under constraints that differ from those of document classification. The messages are short, informal and often written in several languages. Senders alter spelling, insert spaces or replace letters with visually similar symbols in order to evade filters \cite{salman2024,boucher2022}. A filter that runs on a handset or at an operator gateway must classify every message within milliseconds, and the content of private messages should not leave the device or the operator network.

Recent detectors rely more and more on fine-tuned transformer encoders \cite{ghourabi2023,shim2024} and on large language models, and many studies report high accuracy on public corpora of short messages \cite{mishra2023dsmishsms,mishra2022detector,sonowal2020,roy2020,mahmud2024}. Three aspects of the usual evaluation make it hard to tell whether this added capacity is needed in practice.

\emph{Duplicates and leakage.} The Mendeley SMS Phishing Dataset \cite{mendeley2022,mishra2023dataset} is one of the few public corpora that separates smishing from spam. Of its \val{men-raw} messages, our deduplication removed \val{men-removed}: exact or near duplicates of a kept message and members of \val{men-conflict-groups} groups of identical or nearly identical texts that carry different labels. Of the \val{men-clean} remaining messages, \val{men-in-uci} (\val{men-in-uci-pct}\%) have an exact or near duplicate in the older UCI SMS Spam Collection \cite{almeida2011}. A random split of such data can place copies of one message in both the training and the test part, a known source of leakage and of optimistic estimates \cite{kapoor2023,lee2022}. The published methods that we re-implement do not state whether duplicates were removed \cite{mahmud2024,shim2024}.

\emph{Obfuscation that the model has not seen.} Character-level perturbations degrade word-level and subword neural models \cite{gao2018,li2019,pruthi2019}, and Unicode look-alike characters can change the input of a model without changing what a reader sees \cite{boucher2022}. When robustness is measured at all, the test perturbations are usually of the same kind as those used in training, which says little about an attacker who uses a different trick.

\emph{Inference cost.} Accuracy is rarely reported together with latency, and almost never end to end for a complete detection pipeline. In our measurements on one CPU thread, a fine-tuned DistilBERT classifier needs a median of \val{lat-distilaug-p50}~ms per message, while a linear model over character n-grams needs \val{lat-svmchar-p50}~ms.

This paper asks whether transformers are necessary for real-time smishing detection once these three aspects are controlled. We evaluate detectors on three public corpora in English and Portuguese after removing exact and near duplicates, model obfuscation as a random string transformation with one family of operations for training and another family held out for testing, and measure latency on one CPU thread. The detector we study has two operating points. \emph{Lite} is a linear support vector machine over character n-grams with calibrated class probabilities, small and fast enough to run on a handset. \emph{Full} is a cascade for an operator gateway: it answers with the character model when its calibrated confidence reaches a threshold and passes the remaining messages to a DistilBERT classifier. Both are compared with classical, neural and transformer baselines, with a few-shot large language model \cite{qwen2024}, and with re-implementations of two recent published methods \cite{mahmud2024,shim2024} in their own task settings, under a protocol that was frozen before the test data were read.

The contributions are as follows.
\begin{enumerate}
\item \emph{A leakage-free, obfuscation-aware evaluation protocol.} It uses deduplicated splits of three corpora with conflicting labels resolved and no exact or near duplicates across the splits; one obfuscation family for training and a different one held out for testing; a configuration frozen before testing; \val{n-seeds} training seeds; paired bootstrap and McNemar tests with Holm correction; and re-implementation of published methods in the task settings of their papers. We also show that the treatment of duplicates alone changes cross-validated macro F1 by up to \val{infl-max-drop} points, more than most differences between models.
\item \emph{Evidence that published transformer detectors are fragile under unseen obfuscation.} In the smishing-versus-spam setting of Shim et al. \cite{shim2024}, the re-implemented BERT, DistilBERT and RoBERTa detectors reach between \val{sh-clean-min} and \val{sh-clean-max} macro F1 on clean test data and fall to between \val{sh-ho20-min} and \val{sh-ho20-max} when \val{lvl-c}\% of the characters are edited with the held-out operations. Lite goes from \val{svmchar-shim-clean} to \val{svmchar-shim-ho20} on the same data, and its advantage at the \val{lvl-b}\% and \val{lvl-c}\% levels is significant by both tests. Lite needs \val{lat-svmchar-p50}~ms per message, \val{speed-ratio-min} to \val{speed-ratio-max} times less than these transformers.
\item \emph{Two operating points with an explicit choice of the threshold.} The threshold of Full is the solution of a stated selection problem on validation data. In the three-class task, Full sends \val{lat-cascade-routed-pct}\% of the test messages to DistilBERT, has a median end-to-end latency of \val{lat-cascade-p50}~ms and the lowest error cost of all models under the strongest held-out obfuscation. On clean data, Full does not differ significantly from Lite.
\item \emph{Interpretation and deployment.} We read the n-gram weights of the trained Lite model to show which character patterns drive its decisions, present a case study for Kazakhstan with messages in Russian and Kazakh, and derive practical recommendations for banks and mobile operators.
\end{enumerate}

Section~\ref{sec:related} reviews related work. Section~\ref{sec:threat} states the threat model. Section~\ref{sec:methods} describes the data, the obfuscation model, Lite and Full, the cost function, the statistical tests and the protocol. Section~\ref{sec:results} presents the results, Section~\ref{sec:discussion} discusses them, and Section~\ref{sec:conclusions} concludes.
